\section{Circuit Complexity --- Size Classes}
%%% generated by the blueprint pipeline from TCSlib.Complexity.CircuitComplexity.SizeClasses
\subsection*{Circuit Complexity --- Size Classes: Introduction}

Declarations of the Lean module \texttt{TCSlib.Complexity.CircuitComplexity.SizeClasses}.

\subsection*{Circuit Complexity --- Size Classes: Results}

\begin{theorem}[Monotonicity of layered circuit size]
\lean{Language.InLayeredSIZE.mono}
\label{Language.InLayeredSIZE.mono}
\leanok
\uses{Language.InLayeredSIZE}
\difficulty{3}
Let $T,T' : \mathbb{N} \to \mathbb{N}$ and let $L$ be a language of Boolean strings. If
$L$ is recognized by a family of finite-layered circuits using only identity, NOT, and
unbounded fan-in AND gates, with size at most $T(n)$ on inputs of length $n$, and $T(n)
\le T'(n)$ for every $n$, then $L$ is recognized by such a family with size at most
$T'(n)$ on inputs of length $n$.
\end{theorem}

\begin{theorem}[A size bound gives layered polynomial size]
\lean{Language.InLayeredSIZE.inLayeredPPoly}
\label{Language.InLayeredSIZE.inLayeredPPoly}
\leanok
\uses{Language.InLayeredPPoly, Language.InLayeredSIZE, Language.InLayeredSIZE.mono}
\difficulty{1}
Let $T$ be a function from natural numbers to natural numbers, and let $L$ be a language
over Boolean strings. If $L$ is recognized by a finite layered-circuit family using only
the standard gate operations and having size at most $T(n)$ at every input length $n$,
and if there are natural numbers $a,k$ with $T(n)\leq a(n+1)^k$ for every $n$, then $L$
is recognized by such a layered-circuit family with size at most $a(n+1)^k$ at every
input length $n$.
\end{theorem}

\begin{theorem}[Inverse Boolean encoding of one]
\lean{finTwoEquiv_symm_eq_one_iff}
\label{finTwoEquiv_symm_eq_one_iff}
\leanok
\difficulty{4}
For a Boolean value $b$, its image under the inverse of the standard bijection between
Booleans and the two-element type equals $1$ if and only if $b$ is true.
\end{theorem}

\begin{definition}[Language of all-one words]
\lean{Language.allOnes}
\label{Language.allOnes}
\leanok
The language of all-one Boolean words is the set of finite Boolean words in which every
entry is true.
\end{definition}

\begin{theorem}[Membership in the all-ones language]
\lean{Language.mem_allOnes_iff}
\label{Language.mem_allOnes_iff}
\leanok
\uses{Language.allOnes}
\difficulty{1}
A Boolean word $w$ belongs to the language of all-ones words if and only if $w$ is the
list of $w$'s length many true bits.
\end{theorem}

\begin{definition}[Boolean conjunction gate]
\lean{BoolCircuit.andGateOp}
\label{BoolCircuit.andGateOp}
\leanok
\uses{BoolCircuit.GateOp}
For each natural number $w$, define a gate operation with $w$ input slots and alphabet
$	ext{Fin }2$ whose output is the product of the input values in $	ext{Fin }2$
(equivalently, their conjunction).
\end{definition}

\begin{theorem}[AND operations are standard gates]
\lean{BoolCircuit.andGateOp_mem_stdGateOps}
\label{BoolCircuit.andGateOp_mem_stdGateOps}
\leanok
\uses{BoolCircuit.andGateOp, BoolCircuit.stdGateOps}
\difficulty{1}
For every natural number $w$, the operation on Boolean values that takes $w$ inputs and
returns their product belongs to the standard gate set consisting of the identity, NOT,
and all finite-arity AND operations.
\end{theorem}

\begin{theorem}[When a product of bits equals one]
\lean{BoolCircuit.prod_fin_two_eq_one_iff}
\label{BoolCircuit.prod_fin_two_eq_one_iff}
\leanok
\difficulty{4}
Let $\iota$ be a finite index set, and assign to each $i\in\iota$ an element $x_i$ of
the two-element set with multiplication modulo $2$. The product $\prod_{i\in\iota}x_i$
equals $1$ if and only if every $x_i$ equals $1$.
\end{theorem}

\begin{definition}[Nodes in the all-ones circuit]
\lean{BoolCircuit.allOnesNodes}
\label{BoolCircuit.allOnesNodes}
\leanok
For each natural number $n$, define the node types in the two layers of a circuit so
that layer $0$ has type $	ext{Fin }n$ and layer $1$ has type $	ext{Unit}$.
\end{definition}

\begin{definition}[All-ones circuit]
\lean{BoolCircuit.allOnesCircuit}
\label{BoolCircuit.allOnesCircuit}
\leanok
\uses{BoolCircuit.LayeredCircuit, BoolCircuit.allOnesNodes, BoolCircuit.andGateOp}
For each natural number $n$, define a depth-one layered circuit over $	ext{Fin }2$, with
inputs in $	ext{Fin }n$ and a single output. Its input layer has type $	ext{Fin }n$, its
output layer has type $	ext{Unit}$, and its sole output gate takes all $n$ input values
and returns their product in $	ext{Fin }2$.
\end{definition}

\begin{theorem}[Evaluation of the all-ones circuit]
\lean{BoolCircuit.allOnesCircuit_eval₁}
\leanok
\uses{BoolCircuit.LayeredCircuit.eval₁, BoolCircuit.allOnesCircuit}
\difficulty{1}
For every natural number $n$ and every assignment $x :
\mathrm{Fin}(n)\to\mathrm{Fin}(2)$, the all-ones circuit evaluates to
$\prod_{i\in\mathrm{Fin}(n)}x(i)$ in $\mathrm{Fin}(2)$.
\end{theorem}

\begin{theorem}[Size of the all-ones circuit]
\lean{BoolCircuit.allOnesCircuit_size}
\label{BoolCircuit.allOnesCircuit_size}
\leanok
\uses{BoolCircuit.LayeredCircuit.size, BoolCircuit.allOnesCircuit}
\difficulty{1}
For every natural number $n$, the all-ones circuit on $n$ inputs has size $1$.
\end{theorem}

\begin{definition}[Family of all-ones circuits]
\lean{BoolCircuit.allOnesFamily}
\label{BoolCircuit.allOnesFamily}
\leanok
\uses{BoolCircuit.LayeredCircuitFamily, BoolCircuit.allOnesCircuit}
Define a family of circuits indexed by the natural number $n$, where the circuit for
inputs in $	ext{Fin }n$ is the depth-one circuit that outputs the product of all input
values in $	ext{Fin }2$. Every layer of every circuit in this family has a finite node
type.
\end{definition}

\begin{theorem}[Circuit of the all-ones family]
\lean{BoolCircuit.allOnesFamily_circuit}
\label{BoolCircuit.allOnesFamily_circuit}
\leanok
\uses{BoolCircuit.allOnesCircuit, BoolCircuit.allOnesFamily}
\difficulty{1}
For every natural number $n$, the circuit assigned by the all-ones family to inputs of
length $n$ is exactly the depth-one circuit with input nodes indexed by
$Fin\,n$, one output node, and an AND gate taking all $n$ input values.
\end{theorem}

\begin{theorem}[Gate Set of the All-Ones Family]
\lean{BoolCircuit.allOnesFamily_onlyUsesGates}
\label{BoolCircuit.allOnesFamily_onlyUsesGates}
\leanok
\uses{BoolCircuit.LayeredCircuit.onlyUsesGates, BoolCircuit.LayeredCircuitFamily.OnlyUsesGates, BoolCircuit.allOnesCircuit, BoolCircuit.allOnesFamily, BoolCircuit.andGateOp_mem_stdGateOps, BoolCircuit.stdGateOps}
\difficulty{4}
For every natural-number input length, the corresponding circuit in the all-ones family
uses only the identity operation, the negation operation $x \mapsto 1-x_0$, or an AND
operation of some finite arity.
\end{theorem}

\begin{theorem}[Language of the all-ones family]
\lean{BoolCircuit.allOnesFamily_language}
\label{BoolCircuit.allOnesFamily_language}
\leanok
\uses{BoolCircuit.LayeredCircuitFamily.language, BoolCircuit.LayeredCircuitFamily.mem_language_iff, BoolCircuit.allOnesCircuit_eval₁, BoolCircuit.allOnesFamily, BoolCircuit.allOnesFamily_circuit, BoolCircuit.prod_fin_two_eq_one_iff, Language.allOnes, finTwoEquiv_symm_eq_one_iff}
\difficulty{3}
The language recognized by the all-ones circuit family consists exactly of the Boolean
words whose entries are all true.
\end{theorem}

\begin{theorem}[Constant layered size for the all-ones language]
\lean{Language.allOnes_inLayeredSIZE_one}
\label{Language.allOnes_inLayeredSIZE_one}
\leanok
\uses{BoolCircuit.allOnesCircuit_size, BoolCircuit.allOnesFamily, BoolCircuit.allOnesFamily_language, BoolCircuit.allOnesFamily_onlyUsesGates, Language.InLayeredSIZE, Language.allOnes}
\difficulty{1}
The language consisting of Boolean words whose entries are all true is recognized by a
family of finite-layer Boolean feedforward circuits using only identity, NOT, and
unbounded fan-in AND gates, with size at most $1$ for every input length.
\end{theorem}

\begin{theorem}[Linear layered size for the all-ones language]
\lean{Language.allOnes_inLayeredSIZE_linear}
\label{Language.allOnes_inLayeredSIZE_linear}
\leanok
\uses{Language.InLayeredSIZE, Language.InLayeredSIZE.mono, Language.allOnes, Language.allOnes_inLayeredSIZE_one}
\difficulty{1}
The language consisting of Boolean words whose entries are all true is recognized by a
family of finite-layer Boolean feedforward circuits using only identity, NOT, and
unbounded fan-in AND gates, with size at most $n+1$ on inputs of length $n$.
\end{theorem}

\begin{theorem}[Polynomial layered size for the all-ones language]
\lean{Language.allOnes_inLayeredPPoly}
\label{Language.allOnes_inLayeredPPoly}
\leanok
\uses{Language.InLayeredPPoly, Language.InLayeredSIZE.inLayeredPPoly, Language.allOnes, Language.allOnes_inLayeredSIZE_linear}
\difficulty{1}
The language consisting of Boolean words whose entries are all true is recognized by a
family of finite-layer Boolean feedforward circuits using only identity, NOT, and
unbounded fan-in AND gates, with size bounded by a polynomial in the input length.
\end{theorem}
